\documentclass[10pt]{article}
\usepackage[margin=0.62in]{geometry}
\usepackage{amsmath,amssymb}
\usepackage[T1]{fontenc}
\usepackage[hidelinks]{hyperref}
\usepackage{microtype}

\newcommand{\M}{\mathcal M}
\newcommand{\Diff}{\operatorname{Diff}}
\newcommand{\Isom}{\operatorname{Isom}}

\title{Geometric-Scattering Stability to Diffeomorphisms\\
Was Proved in the Expanded Paper}
\author{fa\_banach\_001 / agent\_lane\_00}
\date{17 August 2026}

\begin{document}
\maketitle
\vspace{-1.2em}
\small

\noindent\textbf{Status: literature already answered (qualified
affirmative answer; no new theorem is claimed here).}

\section*{The 2018 target}

Perlmutter, Wolf, and Hirn introduce geometric scattering on compact smooth
Riemannian manifolds.  After proving commutator estimates for its spectral
filters, their introduction says these results should permit a future study
of stability of the \emph{scattering network} to diffeomorphisms, using a
notion analogous to Euclidean Lipschitz stability
\cite{PWH} (source PDF page~3).

This is more than generic future-work language: it identifies the missing
upgrade from individual-filter control to stability of the nonlinear
cascade.

\section*{Explicit later answer}

The authors' expanded paper with Gao, arXiv:1905.10448, carries out that
upgrade \cite{PGWH}.  Let $S_J$ be the windowed geometric wavelet scattering
transform on a compact $d$-manifold $\M$, let $V_\zeta f=f\circ\zeta^{-1}$,
and suppose $f$ is $\lambda$-bandlimited.  Theorem~4.1 (supporting PDF
pages~12--13) proves that, when a diffeomorphism is factored into an
isometric part $\zeta_1$ and a residual part $\zeta_2$,
\[
 \|S_Jf-S_JV_\zeta f\|_{2,2}
 \le C(\M)\!\left[
 2^{-dJ}\|\zeta_1\|_\infty\|U_Jf\|_2
 +\lambda^d\|\zeta_2\|_\infty\|f\|_2\right].
\]
It also gives the finite-depth version, with the first term multiplied by
$(L+1)^{1/2}$ and $\|f\|_2$ factoring the entire bound.

Theorem~4.2 (supporting PDF page~13) sends the averaging scale to infinity
and obtains the non-windowed estimate
\[
 \|\overline Sf-\overline S V_\zeta f\|_2
 \le \lambda^d\|\zeta_2\|_\infty\|f\|_2.
\]
Thus the specific network-level diffeomorphism-stability program from 2018
received a quantitative affirmative answer.

\section*{Idea of the later proof}

Bandlimiting makes deformation of the input controllable:
\[
       \|f-V_\zeta f\|_2
       \le C(\M)\lambda^d\|\zeta\|_\infty\|f\|_2.
\]
The later authors split the deformation into an isometry and a residual
diffeomorphism.  Local isometry invariance controls the first part, while
nonexpansiveness of the scattering cascade propagates the input estimate for
the second.  Combining the two estimates yields Theorem~4.1; taking
$J\to\infty$ removes the window and yields Theorem~4.2.

\section*{Scope of the answer}

The answer is qualified.  Theorems~4.1--4.2 assume bandlimited input.  The
later Section~4.2 obtains an unbandlimited result for \emph{finite-width}
networks on two-point homogeneous manifolds, but explicitly leaves the
general infinite-width problem open.  It also leaves future comparison of
scattering coefficients on two manifolds that are diffeomorphic but not
isometric.  These limitations do not undo the affirmative answer to the
2018 program; they delimit it precisely.

\scriptsize
\begin{thebibliography}{9}
\bibitem{PWH}
M.~Perlmutter, G.~Wolf, and M.~Hirn,
\emph{Geometric Scattering on Manifolds}, arXiv:1812.06968 (2018).

\bibitem{PGWH}
M.~Perlmutter, F.~Gao, G.~Wolf, and M.~Hirn,
\emph{Geometric Wavelet Scattering Networks on Compact Riemannian
Manifolds}, arXiv:1905.10448; Proceedings of Machine Learning Research
\textbf{145} (2021), 570--604.
\end{thebibliography}

\end{document}
